% ===== PRÉAMBULE CANONIQUE — à coller VERBATIM en tête de CHAQUE .tex =====
\documentclass[11pt,a4paper]{article}
\usepackage[utf8]{inputenc}\usepackage[T1]{fontenc}\usepackage{lmodern}
\usepackage[french]{babel}
\usepackage{amsmath,amssymb,mathtools}\usepackage{icomma}
\usepackage[version=4]{mhchem}          % \ce{...} : équations & formules
\usepackage{siunitx}\sisetup{locale=FR} % \qty{}{}, \si{}, \ang{}
\usepackage{chemfig}                    % \chemfig{...} : structures développées
\usepackage{xcolor}\usepackage{colortbl}
\usepackage{tikz}\usepackage{pgfplots}\pgfplotsset{compat=1.18}
\usepackage[most]{tcolorbox}\usepackage{enumitem}\usepackage{array,booktabs}
\usepackage[a4paper,margin=2.2cm]{geometry}
\usetikzlibrary{arrows.meta,calc,angles,quotes,positioning,patterns,backgrounds,shapes.geometric}
\definecolor{primary}{RGB}{222,49,99}\definecolor{primarydark}{RGB}{150,22,55}
\definecolor{defcol}{RGB}{0,114,178}\definecolor{methcol}{RGB}{52,92,148}
\definecolor{excol}{RGB}{0,135,90}\definecolor{attncol}{RGB}{200,40,55}
\tcbset{boxbase/.style={enhanced,breakable,boxrule=0.9pt,arc=2.5pt,
  left=3mm,right=3mm,top=2.3mm,bottom=2.3mm,fonttitle=\bfseries,coltitle=white}}
% --- boîtes TUTORIEL (noms EXACTS, ne pas renommer) ---
\newtcolorbox{cle}[1][]{boxbase,colback=defcol!6,colframe=defcol,
  title={\textsc{À retenir}\ifx\relax#1\relax\relax\else\textmd{ : \itshape #1}\fi}}
\newtcolorbox{methode}[1][]{boxbase,colback=methcol!7,colframe=methcol,sharp corners,
  title={\textsc{Méthode}\ifx\relax#1\relax\relax\else\textmd{ --- \itshape #1}\fi}}
\newtcolorbox{exemplebox}[1][]{boxbase,colback=excol!5,colframe=excol,
  title={\textsc{Exemple}\ifx\relax#1\relax\relax\else\textmd{ : \itshape #1}\fi}}
\newtcolorbox{piege}[1][]{boxbase,colback=attncol!6,colframe=attncol,
  title={\textsc{Piège}\ifx\relax#1\relax\relax\else\textmd{ : \itshape #1}\fi}}
\newtcolorbox{remarque}[1][]{boxbase,colback=black!4,colframe=black!45,
  title={\textsc{Remarque}\ifx\relax#1\relax\relax\else\textmd{ : \itshape #1}\fi}}
% --- boîtes EXERCICE / CORRIGÉ (noms EXACTS, auto-comptées) ---
\newtcolorbox[auto counter]{exo}[1][]{boxbase,colback=primary!4,colframe=primary,
  title={\textsc{Exercice}~\thetcbcounter\ifx\relax#1\relax\relax\else\textmd{ --- \itshape #1}\fi}}
\newtcolorbox[auto counter]{corrige}[1][]{boxbase,colback=excol!4,colframe=excol,
  title={\textsc{Corrigé}~\thetcbcounter\ifx\relax#1\relax\relax\else\textmd{ --- \itshape #1}\fi}}
\newcommand{\R}{\mathbb{R}}\newcommand{\N}{\mathbb{N}}\newcommand{\Z}{\mathbb{Z}}\newcommand{\ds}{\displaystyle}
% (un \usepackage{fancyhdr}+en-tête est possible ; il est ignoré côté web)
% =========================================================================
\begin{document}
\sisetup{per-mode=symbol}
\begin{corrige}[dioxyde d'azote dans l'air pollué]
On convertit puis on divise ; la division fixe le nombre de CS au minimum des données.
\[ [\ce{NO2}]=\frac{0,0361\times10^{-3}\ \si{\gram}}{12,0\times10^{-3}\ \si{\cubic\metre}}
   =0,00300833\ldots\ \si{\gram\per\cubic\metre}. \]
Les deux données ont $3$ CS, donc le résultat aussi (chiffre supprimé $8\ge5$, par excès) :
\[ \boxed{[\ce{NO2}]\approx3,01\times10^{-3}\ \si{\gram\per\cubic\metre}}. \]
\end{corrige}

\begin{corrige}[solubilité du chlorure de sodium]
Règle de trois : $x=\dfrac{36\times93,233}{100}=\dfrac{3356,388}{100}=\qty{33,56388}{\gram}$. En
kilogrammes et à $4$ CS :
\[ \boxed{x\approx\qty{0,03356}{\kilo\gram}=3,356\times10^{-2}\ \si{\kilo\gram}}. \]
\textbf{Commentaire.} L'énoncé impose $4$ CS, on s'y plie ; mais en toute rigueur $\qty{36}{\gram}$ n'a
que $2$ CS, ce qui limiterait le résultat à $2$ CS ($\approx\qty{34}{\gram}$). Il faudrait connaître la
solubilité comme $\qty{36,00}{\gram}$ pour justifier $4$ CS.
\end{corrige}

\begin{corrige}[préparation d'une solution de \ce{NaCl}]
On calcule sans arrondir, on n'arrondit qu'à la fin de chaque étape.
\begin{enumerate}[label=\alph*)]
  \item $n=\dfrac{m}{M}=\dfrac{5,00}{58,44}=0,085558\ldots$ La donnée la plus pauvre, $\qty{5,00}{\gram}$,
        a $3$ CS $\Rightarrow \boxed{n\approx\qty{0,0856}{\mole}}$.
  \item Avec $V=\qty{250,0}{\milli\litre}=\qty{0,2500}{\litre}$ :
        $c=\dfrac{n}{V}=\dfrac{0,085558}{0,2500}=0,34223\ldots$ Minimum de CS $=3$ ($5,00$) $\Rightarrow
        \boxed{c\approx\qty{0,342}{\mol\per\litre}}$.
\end{enumerate}
\end{corrige}

\begin{corrige}[acidité de deux solutions]
\begin{enumerate}[label=\alph*)]
  \item $\mathrm{pH}=-\log(4,0\times10^{-2})=1,39794\ldots$ ; $[\ce{H3O+}]$ a $2$ CS $\Rightarrow 2$
        décimales $\Rightarrow \boxed{\mathrm{pH}\approx1,40}$.
  \item $[\ce{H3O+}]=10^{-5,00}=\boxed{1,0\times10^{-5}\ \si{\mol\per\litre}}$ : les $2$ décimales du pH
        donnent $2$ CS.
\end{enumerate}
\end{corrige}

\end{document}
